\documentclass{rapportECL2024}

\addbibresource{biblio.bib}

%--------------------------------------

\titre{Titre}
\soustitre{Sous-titre}

\enseignant{Prénom \textsc{Nom} \\
            Prénom \textsc{Nom} }

\eleves{Prénom \textsc{Nom} \\
	Prénom \textsc{Nom} \\ 
	Prénom \textsc{Nom} }

%--------------------------------------

\begin{document}

\fairepagedegarde
\fairetabledesmatieres

%--------------------------------------

\section{Section 1}
\lipsum[1]

\subsection{Sous-section 1}
\lipsum[2]

\subsection{Sous-section 2}
\lipsum[3]

\subsubsection{Sous-sous-section 1}
\lipsum[4]

\section{Section 2}
\lipsum[5]

\section{Quelques commandes utiles}

\textbf{Insérer une image}
\begin{figure}[H]
    \centering
    \includegraphics[width=0.5\textwidth]{logos/logo.png}
    \caption{Insérer une image.}
    \label{fig:logo}
\end{figure}

ou la version raccourcie avec \texttt{\textbackslash insererfigure} : \insererfigure{logos/logo.png}{0.5\textwidth}{Insérer une image avec une commande.}{logo_raccourci}


\textbf{Citer une figure} avec \texttt{\textbackslash ref} : \ref{fig:logo_raccourci}.\\


\textbf{Insérer des expressions mathématiques} en bloc :
\begin{equation*}
    \int_{-\infty}^{\infty} e^{-x^2} \, \text{d}x = \sqrt{\pi}
\end{equation*}

ou bien en \textit{inline} : $\int_{-\infty}^{\infty} e^{-x^2} \, \text{d}x = \sqrt{\pi}$.\\


\textbf{Insérer une série de calculs :}
\begin{align*}
    x &= 0.999\ldots \\
    10x &= 9.999\ldots \\
    10x - x &= 9.999\ldots - 0.999\ldots \\
    9x &= 9 \\
    x &= 1 \\
    0.999\ldots &= 1
\end{align*}


\textbf{Insérer une liste :}
\begin{itemize}
    \item Premier niveau
    \begin{itemize}
        \item Deuxième niveau
        \item Un autre élément au deuxième niveau
    \end{itemize}
    \item Un autre élément au premier niveau
\end{itemize}


\textbf{Insérer un tableau simple} \\
\begin{table}[H]
    \centering
    \begin{tabular}{|c|c|c|}
        \hline
        A & B & C \\
        \hline
        1 & 2 & 3 \\
        4 & 5 & 6 \\
        \hline
    \end{tabular}
    \caption{Un tableau simple.}
    \label{tab:tableau_simple}
\end{table}


\textbf{Insérer du code en \textit{inline}} : \texttt{print("Hello, World!")}.\\


\textbf{Insérer et référencer une boîte de code} : Pour l'environnement \texttt{codebox}, il faut d'abord renseigner le label de la boîte, suivi du titre, puis du langage dans lequel est écrit le code (la liste des langages pris en charge est disponible \href{https://www.overleaf.com/learn/latex/Code_Highlighting_with_minted#:~:text=Main%20supported%20programming%20languages%20and%20configuration%20files}{\textcolor{blue}{\underline{ici}}}). 

On peut faire référence au code avec : Code \ref{code:code_example}.

\begin{codebox}[code:code_example]{Un exemple de boîte de code \texttt{Python}}{python}
def sieve_of_eratosthenes(n: int) -> list[int]:
    '''Returns a list of all prime numbers up to n.'''
    primes = [True] * (n + 1)
    primes[0] = primes[1] = False  # 0 and 1 are not prime numbers
    for i in range(2, int(n**0.5) + 1):
        if primes[i]:
            for j in range(i * i, n + 1, i):
                primes[j] = False
    return [i for i, is_prime in enumerate(primes) if is_prime]
\end{codebox}


\textbf{Citer une source depuis le fichier \texttt{biblio.bib}} avec \texttt{\textbackslash cite} : \cite{exemple_de_source}.\\


\textbf{Afficher une bibliographie} rapidement avec \texttt{\textbackslash insererbiblio}.

\insererbiblio 

\end{document}